\chapter*{Interlude: What a Horizon Really Is}
\addcontentsline{toc}{chapter}{Interlude: What a Horizon Really Is}
\markboth{INTERLUDE}{INTERLUDE}

\begin{oberflaeche}
Let us pause for a moment, entirely without formulas.

A horizon is not a length and not a format. It is the space a
number currently lives in -- with a fixed, finite number of places
that come from the factorials: two, six, twenty-four, one hundred
and twenty. An ordinary number knows only its value. A Hori number
knows its value \emph{and} its space, and both belong to it the way
a book has both its text and the shelf it stands on.

That sounds like a small thing, but it is a change of perspective.
Our usual mathematics takes the infinite number line as given and
places all numbers on it at once. Horimetrik builds staircases
instead: each step is finite and completely surveyable, and the
infinity lies not in any single step but in the promise that there
is always a next one. Nothing about classical infinity is refuted
by this -- it is merely told differently: as a sequence of spaces
rather than as one single space.

The rest of this book investigates what this way of telling costs
and what it earns. It costs convenience: with every move from step
to step, all digits change their meaning, and you must decide what
the move is to preserve -- the value or the shape. It earns
structure: out of exactly this decision arise the two identities,
the two worlds of arithmetic, and the theorem that you cannot live
in both at once.
\end{oberflaeche}
